\section{Modelos generativos}

A cada imagen generada con un NCA \cite{mordvintsevNCA} le corresponde una ANN que está explícitamente entrenada para imitar ese comportamiento. Alejandro Hernández et al. \cite{NCAM} proponen una ANN de mayor capacidad que permita crecer múltiples imágenes, a lo que llamaron \textit{Neural Cellular Automata Manifold} (NCAM).

El NCAM usa un modelo de convoluciones dinámicas para modelar un espacio de NCAs. Utiliza un \textit{autoencoder} que se compone de un codificador y un decodificador: El codificador transforma la entrada en una representación continua y luego categórica, que el decodificador reconstruye en el espacio original.

Un elemento destacable de la arquitectura propuesta es un parámetro entrenable llamado \textit{Leak Factor} (LF) que mejora la estabilidad durante el entrenamiento. Este parámetro regula la cantidad de información que se filtra en cada paso del modelo, por ejemplo, un valor bajo retiene el conocimiento acumulado. Este valor se incrementa durante el entrenamiento a medida que mejora el rendimiento.

La arquitectura también emplea convoluciones dinámicas. Los \textit{kernels} se generan a partir del \textit{embedding} (incrustación) de cada instancia, permitiendo adaptabilidad sin modificar la estructura del NCA. El NCAM inicia desde una imagen \textit{seed} y actualiza cada celda mediante percepciones locales (filtros Sobel $+$ normalización) y capas densas por píxel, moduladas por el LF. Se exploran actualizaciones estocásticas y sincrónicas, siendo esta última más precisa y estable, tal como en \cite{NiklassonNCA}.

Los experimentos confirman que la codificación y la actualización sincrónica mejoran la calidad de reconstrucción, incluso en ausencia de mecanismos de enmascaramiento. También se observa que aumentar la dimensionalidad del \textit{embedding} mejora significativamente los resultados.

Por otro lado, Kevin Frans en \cite{StampCA} propone un modelo llamado \textit{Conditional Neural Cellular Automata} (StampCA), cuya \textit{seed} recibe un \textit{embedding}, el cual define la imagen en la cual va a crecer. 

Dado un diseño deseado, una red de codificador lo comprime en un \textit{embedding}, el cual se utiliza como condición inicial en la NCA, que luego itera para crecer en la figura objetivo; esto permite crear nuevas configuraciones.

Utilizando un \textit{dataset} (conjunto de datos) con aproximadamente 3000 \textit{emojis} de $32 \times 32$ píxeles, fue capaz de entrenar a StampCA utilizando el mismo método de \cite{mordvintsevNCA} para poder dibujar todos los \textit{emojis} con una única NCA.

Para generar números de MNIST, se adopta un enfoque tipo \textit{Generative Adversarial Network} (GAN), donde el objetivo es crecer dígitos desde ruido aleatorio. Dado que entrenar directamente un NCA como generador GAN resultó inestable, se utilizó una estrategia en dos etapas: (1) Se entrena primero un generador tradicional (\textit{feedforward}) para mapear vectores aleatorios a dígitos MNIST y (2) el NCA se entrena para imitar la salida de ese generador, logrando una generación más estable.

Este enfoque permite que el NCA genere dígitos sin necesidad de un \textit{dataset} explícito, aprendiendo directamente a partir del objetivo GAN. Además, se observa un patrón de crecimiento distintivo: mientras los \textit{emojis} crecen en ondas desde el centro, los dígitos MNIST se desarrollan siguiendo la forma del trazo.

Mingxiang Chen et al. \cite{VAENCA} combinan los principios de \textit{Variational Auto-Encoder} (VAE) con NCA. Utilizan un \textit{encoder} que transforma una imagen en una distribución de probabilidad, las cuales forman un \textit{embedding}, el cual funciona como condición inicial del NCA. Posteriormente, Rasmus Berg Palm et al. \cite{VNCA} nombran a este modelo como \textit{Variational Neural Cellular Automata} (VNCA), quienes además aseguran que ninguno de los otros modelos vistos (\cite{mordvintsevNCA}, \cite{NCAM}, \cite{GCA}, \cite{horibe} y \cite{StampCA}) es realmente generativo.

El VNCA se basa en la arquitectura del modelo generativo probabilístico \textit{Variational Auto-Encoder} (VAE) la cual aprende una distribución de probabilidad que codifica los datos observados. A partir de esta representación, el decodificador genera imágenes mediante un NCA que simula el crecimiento espacial desde una \textit{seed}. A diferencia de VAE, que decodifica la imagen directamente, el VNCA utiliza el \textit{embedding} como configuración inicial de una cuadrícula de celdas que evolucionan siguiendo las reglas definidas por una CNN. 

Por otra parte, este modelo incorpora mecanismos de duplicación celular inspirados en la mitosis. Cada $M$ pasos, las celdas se duplican hacia la derecha y hacia abajo, permitiendo a la imagen crecer de manera autoorganizada. 

Los experimentos realizados en \cite{VNCA} demuestran que el VNCA es capaz de aprender procesos generativos autoorganizados tanto en dominios simples como MNIST, como en conjuntos más complejos como CelebA\footnote{\textit{Dataset} de imágenes de celebridades}. En MNIST, el modelo logra reconstrucciones precisas y muestras plausibles desde el espacio latente, con una log-verosimilitud de $\log p(x) \geq -84.23$ nats\footnote{Unidad de medida utilizada en aprendizaje automático y estadística para expresar la $\log$-verosimilitud o la entropía de un modelo}, su rendimiento es comparable a modelos autorregresivos como PixelCNN. Además, el análisis del espacio latente revela una estructura más separable y organizada que la de decodificadores convolucionales tradicionales, lo que sugiere una codificación más semántica y eficiente.

En CelebA, aunque las muestras generadas directamente no alcanzan calidad visual comparable a modelos de última generación, el VNCA muestra capacidad de reconstrucción robusta y mejora visual significativa al ajustar el parámetro $\beta$, que regula el equilibrio entre fidelidad de reconstrucción y calidad de generación. 

Finalmente, se explora la resiliencia del VNCA ante daño estructural durante el crecimiento. Al eliminar las capas de duplicación y entrenar con un esquema de recuperación de daño, el modelo aprende dinámicas estables capaces de restaurar imágenes deterioradas casi por completo, incluso tras múltiples pasos de evolución. Este comportamiento sugiere que el VNCA no solo codifica la imagen, sino que también es un proceso de reconstrucción en el estado celular.
